\section*{Registro delle Modifiche}

\renewcommand{\arraystretch}{1.5}
\begin{longtable}{>{\ttfamily\small}l l l l p{4cm}}
	\toprule
	\textbf{Versione} & \textbf{Data} & \textbf{Autore} & \textbf{Verificatore} & \textbf{Descrizione} \\
	\midrule
	\endfirsthead

	\toprule
	\textbf{Versione} & \textbf{Data} & \textbf{Autore} & \textbf{Verificatore} & \textbf{Descrizione} \\
	\midrule
	\endhead

	\midrule
	\multicolumn{5}{r}{\textit{Continua nella pagina successiva}} \\
	\midrule
	\endfoot

	\bottomrule
	\endlastfoot
	
	0.3.1 & 2026-08-06 & {Nome Cognome} & {Nome Cognome} & Editing di forma su Introduzione (rimosso paragrafo ridondante rispetto allo Scopo del documento; tolto dettaglio non necessario sul Tedesco; riformulata la descrizione del cappello blu) e Tecnologie (corretto lo strumento di modellazione UML in StarUML; aggiunta nota su DI e pattern non forniti nativamente da Vue.js). \\
	0.3.0 & 2026-08-06 & {Nome Cognome} & {Nome Cognome} & Rimossa la persistenza server-side (PostgreSQL) dallo stack tecnologico e dall'architettura di deployment, non implementata in questa fase. Requisiti opzionali correlati (R4-V-Op, R7/R8-Q-Op, R78--R86) riportati con stato ``Non implementato'' e relativa motivazione nella Sezione 7. \\
	0.2.0 & 2026-08-06 & {Nome Cognome} & {Nome Cognome} & Disciplina dei riferimenti: URL puntuali (non generici) verso la documentazione di terze parti con data di ultimo accesso; citazioni interne circoscritte alla sezione rilevante di Analisi dei Requisiti e Analisi dello Stato dell'Arte, come richiesto in sede di revisione. \\
	0.1.0 & 2026-08-06 & {Nome Cognome} & {Nome Cognome} & Prima stesura completa del documento: tecnologie, architettura logica e di deployment, interfaccia API, design pattern, flusso dati e tracciamento dei requisiti, a partire da Analisi dello Stato dell'Arte v1.0.0 e Analisi dei Requisiti v1.8.0. \\
	0.0.0 & YYYY-MM-DD & {Nome Cognome} & {Nome Cognome} & Prima stesura della struttura del documento. \\
\end{longtable}
